% Generated draft: requires mathematical and source review.
\documentclass[11pt]{article}
\usepackage[margin=25mm]{geometry}
\usepackage{fontspec}
\setmainfont{Latin Modern Roman}
\usepackage{amsmath,amssymb,mathtools,unicode-math}
\setmathfont{Latin Modern Math}
\usepackage{enumitem,needspace,xurl,hyperref}
\setlength{\parindent}{0pt}
\setlength{\parskip}{5pt}
\setlength{\emergencystretch}{3em}
\newcommand{\N}{\mathbb{N}}
\newcommand{\Z}{\mathbb{Z}}
\newcommand{\Q}{\mathbb{Q}}
\newcommand{\R}{\mathbb{R}}
\newcommand{\C}{\mathbb{C}}
\newcommand{\F}{\mathbb{F}}
\newcommand{\A}{\mathbb{A}}
\newcommand{\PP}{\mathbb{P}}
\newcommand{\E}{\mathbb{E}}
\newcommand{\eps}{\varepsilon}
\newcommand{\dd}{\,\mathrm{d}}
\begin{document}
\section*{UnsolvedMath ID: 8700057}
\textbf{Source problem code:} catalogue entry
\subsection{Definitions and mathematical statement}
For an algebraic number $\xi$, let $\mathrm h(\xi)$ denote absolute logarithmic Weil height, with $\mathrm h(0)=0$. For nonzero $\xi$ with primitive minimal polynomial $a_d\prod_{i=1}^d(X-\xi_i)$, the convention is
\[
 \mathrm h(\xi)=\frac1d\left(\log|a_d|+
                    \sum_{i=1}^d\log\max\{1,|\xi_i|\}\right).
\]
All logarithms of positive real quantities below are natural logarithms.

The source asserts existence of a positive absolute constant $C$, uniform in the following data. Let $n\ge1$, let $\alpha_1,\ldots,\alpha_n\in\overline\Q^\times$, and choose complex logarithms $\log\alpha_j$ with $\exp(\log\alpha_j)=\alpha_j$. Suppose these chosen logarithms are linearly independent over $\Q$. Let $\beta_0,\ldots,\beta_n\in\overline\Q$ be not all zero, and put
\[
 D=[\Q(\alpha_1,\ldots,\alpha_n,
                  \beta_0,\ldots,\beta_n):\Q].
\]
Choose real $A_1,\ldots,A_n,B\ge e$ such that
\[
 \log A_j\ge\max\{\mathrm h(\alpha_j),
                    |\log\alpha_j|/D,1/D\}\quad(1\le j\le n),
 \qquad
 B\ge\max_{1\le j\le n-1}\mathrm h(\beta_j).
\]
For
\[
 \Lambda=\beta_0+\sum_{j=1}^n\beta_j\log\alpha_j,
\]
the precise extracted lower bound is
\[
 |\Lambda|>
 \exp\left\{-C^nD^{n+2}
       \left(\prod_{j=1}^n\log A_j\right)
       (\log B+\log D)(\log D)\right\}.
\]
The particular height condition and all factors are retained literally.

There are material source-formulation problems. The $B$ condition controls only $\beta_1,\ldots,\beta_{n-1}$, omitting $\beta_0$ and $\beta_n$, and uses $B$ rather than $\log B$ as written. At $n=1$ the displayed maximum is over an empty set. Moreover $D=1$ is not excluded and makes the final $\log D$ factor zero, so the asserted right-hand side becomes $1$. Those boundary and coefficient-control cases require checking against the source; no missing height bound, regularized logarithm, or degree restriction is guessed. The statement also does not separately assume $\Lambda\ne0$, and none is silently added.
\subsection{Short English statement}
The conjecture proposes a specific absolute-constant lower bound for an inhomogeneous algebraic linear form in chosen independent logarithms. Its exact dependence on field degree, heights, and logarithm sizes is the target; existence of some effective logarithmic-form estimate is not enough.

The supplied summary mentions Baker and Matveev as providing related effective estimates but says the exact formula is not established verbatim. It characterizes the issue as quantitative refinement, while supplying no theorem matching these algebraic coefficients and the literal formula's boundary cases. The OPEN-TRIAGE label and 2004 survey are retained with current primary-literature status audit required.
\subsection{Sources}
\begin{itemize}
\item Local UnsolvedMath registry, permanent ID 8700057, source code catalogue entry, ``Conjecture 4.15,'' proposer metadata Michel Waldschmidt. Dataset curation and metadata: Ulam AI and UnsolvedMath Contributors, CC BY 4.0.
\item Michel Waldschmidt, \emph{Open Diophantine Problems} (2004), Conjecture 4.15, \url{https://arxiv.org/abs/math/0312440}; source-TeX extraction, accessed 2026-07-29.
\item Baker and Matveev are named in the supplied research summary without complete references or exact hypotheses supporting this particular inhomogeneous algebraic-coefficient estimate.
\item Source-specific provenance: the survey presents the item as open; current status audit and redistribution terms NEEDS\_REVIEW for the publicly accessible author/publisher copy, separately from the dataset licence.
\end{itemize}
\end{document}
